A crescente adoção de modelos de linguagem de grande escala em sistemas de produção torna os \textit{guardrails}, mecanismos de classificação de segurança que inspecionam entradas e saídas de LLMs, componentes críticos de infraestruturas que atendem usuários em múltiplos idiomas. Em ambientes de implantação local, a quantização pós-treinamento é amplamente adotada para reduzir os requisitos de memória e aumentar o desempenho desses modelos. Contudo, o impacto da quantização sobre o desempenho de \textit{guardrails} em idiomas distintos do inglês, especialmente o português brasileiro, permanece sem investigação sistemática na literatura. Este trabalho avalia como a quantização afeta o \textit{gap} de desempenho entre inglês e português brasileiro em \textit{guardrails} baseados em LLM, operacionalizado pela métrica $\Delta$F1. Cinco modelos são avaliados, Llama Guard 3 8B, Llama Guard 3 1B, ShieldGemma 9B, Nemotron Safety 8B e Granite Guardian 4.1 8B, em três níveis de precisão numérica: FP16, GPTQ 8\textit{bits} e GPTQ 4\textit{bits}, sobre um conjunto de 16.188 prompts adaptados do BeaverTails e traduzidos para PT-BR via Google Translate. O experimento adota Bootstrap com 1.000 iterações para estimação de intervalos de confiança e o teste de McNemar para verificação de significância estatística do \textit{gap} entre idiomas. 
% Separe as palavras-chave por ponto
\palavraschave{\textit{guardrails}. quantização. segurança de LLMs. português brasileiro. multilinguismo.}